% translation-source: docs/user/latex/chapters/07-couleur.tex; sha256: 26daf42a5ca80611cc543e8d13bc3576ca31df4a090b8ca9af27366d4c953719
\chapter{Global color}
\section{Temperature and tint}
Temperature spans 2\,000 to 50\,000 K; tint spans $-150$ to  $+150$.
Temperature shifts warm/cool balance; tint corrects the green/magenta axis.
Slider bands are indications, not a contractual swatch of the result.
Automatic white balance is in the RAW group.

The reference temperature depends on photo development; it is not always
5\,200 K or6\,500 K. Reset returns to that reference, not zero kelvin.
Tint resets to zero in the current interface. Units or values shown by
another application do not guarantee exact numeric correspondence.

\section{Vibrance and saturation}
Both controls span$-100$ to  $+100$, with zero as reference. Saturation changes
overall color richness. Vibrance seeks a more selective effect based on
colors already present; it is not simply weaker saturation. Strong effects
on vivid colors may also change their apparent brightness.

\paragraph{Example.}
Begin with a small temperature correction, then adjust tint. Add moderate
vibrance for dull colors and inspect skin tones, reds and yellows.
Adjust global saturation while viewing both the whole image and already
colorful regions. Hold \touche{Space} in the view to compare with the original.

\attention{Color validations depend on the starting rendition, profile and
camera. Extreme values are useful tests, not a promise to reproduce Lightroom
exactly.}
\gestesreglages
